\begin{longtable}[c]{| L{6cm} | L{6cm} |}
\caption{UART 和 LP UART 特性区分}
\label{tab:uart-feature-distinction} \\
    \hline
    \rowcolor{lightgray}
    \textbf{UART 特性} & \textbf{LP\_UART 特性}\\\hline
    \endfirsthead
    
    \multicolumn{2}{c}{\bfseries \tablename\ \thetable{} -- 接上页}\\\hline
    \rowcolor{lightgray}
    \textbf{UART 特性} & \textbf{LP\_UART 特性}\\\hline
    \endhead
    
    \multicolumn{2}{c}{\textbf{见下页}} \\
    \endfoot    
    \endlastfoot
    
\multicolumn{2}{|c|}{可编程收发波特率，最大为 5 MBaud} \\ \hline
一个 UART 的发送通道和接收通道分别有 128 x 8 位 RAM 存储空间 & LP UART 的发送通道和接收通道分别有 32 x 8 bit RAM 存储空间 \\ \hline
\multicolumn{2}{|c|}{全双工异步通信} \\ \hline
\multicolumn{2}{|c|}{数据位（5 到 8 位）} \\ \hline
\multicolumn{2}{|c|}{停止位 （1、1.5 或 2 位）} \\ \hline
\multicolumn{2}{|c|}{奇偶校验位} \\ \hline
\multicolumn{2}{|c|}{AT\_CMD 特殊字符检测} \\ \hline
RS485 协议 & --- \\ \hline
IrDA 协议 & --- \\ \hline
GDMA 高速数据通信 & --- \\ \hline
\multicolumn{2}{|c|}{接收超时} \\\hline
\multicolumn{2}{|c|}{UART 唤醒模式} \\ \hline
\multicolumn{2}{|c|}{软件流控和硬件流控} \\ \hline
\makecell[l]{三个可预分频的时钟源：\\1. APB\_CLK\\ 2. XTAL\_CLK\\ 3. RC\_FAST\_CLK} & \makecell[l]{两个可预分频的时钟源：\\1. XTAL\_D2\_CLK\\ 2. LP\_FAST\_CLK} \\ \hline
\end{longtable}